%!TEX root = tfg.tex
\chapter{Metodología}

\begin{quotation}[General and President]{Dwight D. Eisenhower (1890--1969)}
Plans are useless, but planning is indispensable.
\end{quotation}

\begin{abstract}
En este capítulo se expone el marco metodológico que ha guiado el desarrollo de
\textit{Via Inferni}. La propuesta combina prácticas de planificación inspiradas en
PMBOK, un enfoque iterativo basado en FDD y una gestión del código orientada a mantener
trazabilidad y estabilidad durante el desarrollo.
\end{abstract}

\section{PMBOK adaptado a un videojuego}
El proyecto toma como referencia varias áreas de conocimiento del PMBOK, pero adaptadas a
la escala real de un TFG y a las particularidades del desarrollo de un videojuego. En
lugar de aplicar una gestión documental pesada, se seleccionan únicamente aquellos planes
que aportan control efectivo sobre el alcance, el calendario, la comunicación, los
riesgos, la calidad y los costes.

Este enfoque permite mantener una visión de proyecto completa sin perder agilidad. La
planificación sirve como marco de decisión, mientras que la ejecución técnica se documenta
de forma incremental y vinculada a entregables funcionales.

\section{Feature Driven Development}
Para la construcción de \textit{Via Inferni} se ha seleccionado la metodología
\textit{Feature Driven Development} (FDD). La razón principal es que encaja especialmente
bien en proyectos donde la funcionalidad debe crecer por incrementos pequeños,
comprobables y con un diseño técnico suficientemente claro antes de implementar cada
cambio.

FDD permite trabajar con una visión global del producto, pero sin perder granularidad en
la ejecución. En lugar de abordar el desarrollo como un bloque único, el proyecto se
descompone en funcionalidades manejables que pueden diseñarse, implementarse y validarse
por separado.

El proceso se articula a través de cinco pasos fundamentales que guían el proyecto desde
su concepción hasta la incorporación de cada funcionalidad, tal y como se ilustra en la
Cuadro \ref{fig:fdd_flow}.

\begin{figure}[htb]
    \centering
    \begin{tikzpicture}[node distance=1.2cm, auto,
        block/.style={rectangle, draw, fill=gray!5, text width=6cm, text centered, rounded corners, minimum height=1cm, font=\small},
        line/.style={draw, -latex, thick}]
        \node [block] (step1) {\textbf{1. Desarrollo de un modelo general} \\ (Análisis de arquitectura y dominio)};
        \node [block, below=of step1] (step2) {\textbf{2. Construcción de la lista de funcionalidades} \\ (Desglose del backlog)};
        \node [block, below=of step2] (step3) {\textbf{3. Planificación por funcionalidad} \\ (Priorización e hitos)};
        \node [block, below=of step3, fill=blue!5] (step4) {\textbf{4. Diseño por funcionalidad} \\ (Contratos y lógica técnica)};
        \node [block, below=of step4, fill=blue!5] (step5) {\textbf{5. Construcción por funcionalidad} \\ (Codificación e integración)};

        \path [line] (step1) -- (step2);
        \path [line] (step2) -- (step3);
        \path [line] (step3) -- (step4);
        \path [line] (step4) -- (step5);
        \path [line] (step5.east) -- ++(1,0) -- ++(0,2.2) -- (step4.east);
        \node [right=of step4, xshift=0.3cm, yshift=1cm, text width=2cm, font=\footnotesize\itshape] {Iteración continua por funcionalidad};
        \draw[dashed, blue!40, thick] ([xshift=-0.5cm, yshift=0.3cm]step4.north west) rectangle ([xshift=2.5cm, yshift=-0.3cm]step5.south east);
        \node[blue!60, rotate=90, anchor=south] at ([xshift=2.5cm]step5.east) {Fase iterativa};
    \end{tikzpicture}
    \caption{Flujo de procesos de la metodología FDD aplicado al desarrollo del sistema.}
    \label{fig:fdd_flow}
\end{figure}

\section{Gestión de la configuración y el código}
La gestión del repositorio es crítica para asegurar la integridad del proyecto. Para ello
se ha definido una estrategia híbrida que combina la organización estructural de
\textit{GitFlow} con la integración frecuente propia de \textit{Trunk-Based Development}.

La rama \texttt{main} representa el estado estable del proyecto, mientras que
\texttt{develop} actúa como punto de integración del trabajo validado. Las ramas de
funcionalidad son efímeras y se crean para resolver tareas concretas del \textit{backlog},
reduciendo así el coste de integración y el riesgo de divergencias largas.

Como criterio de trazabilidad se ha adoptado el estándar \textit{Conventional Commits},
que facilita identificar la intención de cada cambio y relacionarlo con funcionalidades,
correcciones o ajustes de documentación.

\section{Herramientas y tecnologías}
El desarrollo del proyecto se apoya en un conjunto de herramientas técnicas y de gestión
que permiten coordinar el trabajo y mantener control sobre el avance.

\begin{itemize}
    \item Unity actúa como motor principal de desarrollo del videojuego.
    \item Git y GitHub se utilizan para el control de versiones, revisión de cambios y seguimiento de tareas.
    \item GitHub Projects proporciona el tablero Kanban con el que se monitoriza el flujo de trabajo.
    \item La memoria se mantiene en LaTeX, lo que permite modularizar capítulos y gestionar el documento con el mismo rigor que el código.
\end{itemize}
